\documentclass[12pt]{amsart}

\usepackage{amsmath}
\usepackage{amssymb}
\usepackage{amsthm}
\usepackage{mathtools}
\usepackage{enumitem}
\usepackage[colorlinks=true,citecolor=blue,urlcolor=blue,linkcolor=blue]{hyperref}

% Theorem environments
\newtheorem{theorem}{Theorem}[section]
\newtheorem{lemma}[theorem]{Lemma}
\newtheorem{proposition}[theorem]{Proposition}
\newtheorem{corollary}[theorem]{Corollary}
\newtheorem{conjecture}[theorem]{Conjecture}

\theoremstyle{definition}
\newtheorem{definition}[theorem]{Definition}
\newtheorem{assumption}[theorem]{Assumption}
\newtheorem{example}[theorem]{Example}

\theoremstyle{remark}
\newtheorem{remark}[theorem]{Remark}

% Operators
\DeclareMathOperator{\tr}{tr}
\DeclareMathOperator{\diag}{diag}
\DeclareMathOperator{\rank}{rank}
\DeclareMathOperator{\MSE}{MSE}
\DeclareMathOperator{\relu}{ReLU}
\DeclareMathOperator{\spn}{span}
\DeclareMathOperator{\Var}{Var}
\DeclareMathOperator{\Fisher}{Fisher}
\DeclareMathOperator{\softmax}{softmax}

% Shortcuts
\newcommand{\R}{\mathbb{R}}
\newcommand{\E}{\mathbb{E}}
\newcommand{\Prob}{\mathbb{P}}
\newcommand{\bx}{\mathbf{x}}
\newcommand{\bX}{\mathbf{X}}
\newcommand{\bJ}{\mathbf{J}}
\newcommand{\bM}{\mathbf{M}}
\newcommand{\bU}{\mathbf{U}}
\newcommand{\bV}{\mathbf{V}}
\newcommand{\bI}{\mathbf{I}}
\newcommand{\bu}{\mathbf{u}}
\newcommand{\bv}{\mathbf{v}}
\newcommand{\bw}{\mathbf{w}}
\newcommand{\ba}{\mathbf{a}}
\newcommand{\bg}{\mathbf{g}}
\newcommand{\bp}{\mathbf{p}}
\newcommand{\bH}{\mathbf{H}}
\newcommand{\bF}{\mathbf{F}}
\newcommand{\bA}{\mathbf{A}}
\newcommand{\bB}{\mathbf{B}}
\newcommand{\bP}{\mathbf{P}}
\newcommand{\by}{\mathbf{y}}
\newcommand{\bSigma}{\boldsymbol{\Sigma}}
\newcommand{\bGamma}{\boldsymbol{\Gamma}}
\newcommand{\bxi}{\boldsymbol{\xi}}
\newcommand{\bmu}{\boldsymbol{\mu}}
\newcommand{\norm}[1]{\left\lVert #1 \right\rVert}
\newcommand{\abs}[1]{\left\lvert #1 \right\rvert}
\newcommand{\ip}[2]{\langle #1, #2 \rangle}

\title[Cross-Entropy Noise Propagation and Neural Collapse]{%
  A Cross-Entropy Noise Propagation Operator and\\
  the Effective Rank at the Neural Collapse Fixed Point}

\author{Lightman Chang}
\address{Independent Researcher}
\email{lightman.chang@gmail.com}

\date{\today}

\subjclass[2020]{Primary 68T07; Secondary 62G05, 41A25, 60J60, 62H30}

\keywords{Cross-entropy loss, Neural Collapse, Fisher information,
  noise propagation, benign overfitting, effective rank,
  multi-class classification}

\begin{document}

\begin{abstract}
We extend the noise-propagation framework
of~\cite{paperA}---originally formulated for the squared loss in
regression---to multi-class classification under the cross-entropy
loss. The technical step is to derive a Fisher-corrected noise
propagation operator $\bM_{\mathrm{CE}}$ directly from the
cross-entropy Hessian
$\nabla_{\theta}^{2}\ell_{\mathrm{CE}} = \bJ^{\top}\bH\bJ$, where
$\bH(\bx) = \diag(\bp(\bx)) - \bp(\bx)\bp(\bx)^{\top}$ is the
output-space Fisher matrix at the softmax probabilities $\bp(\bx)$.
Our main result (Theorem~\ref{thm:ce-trichotomy}) is that
$\tr(\bM_{\mathrm{CE}}) < \infty$ if and only if the underlying
power-law exponents satisfy $\beta > \alpha + 1/2$, identical to the
trichotomy for MSE. We then study the limit of $\bM_{\mathrm{CE}}$
along a sequence of trained predictors approaching the Neural
Collapse fixed point of Papyan--Han--Donoho. Although
$\bH$ becomes singular in the strict one-hot limit, the limiting
operator that controls the test-risk effective rank is well-defined
on the $(K-1)$-dimensional zero-sum subspace fixed by the simplex
equiangular tight frame, and Theorem~\ref{thm:ce-rank} computes the
effective rank of the per-sample block (equivalently, of the
output-space averaged operator $\overline{\bM}_{\mathrm{CE}}$ on the
zero-sum subspace $\mathcal{Z}_{K}$): it equals exactly $K-1$, the
number of classes minus one. The argument is purely linear-algebraic
given the simplex ETF geometry, and yields a dimension-free
generalization scaling for $K$-class classification at the terminal
phase of training.
\end{abstract}

\maketitle

%% ============================================================
\section{Introduction}\label{sec:intro}
%% ============================================================

The unified bound of~\cite{paperA} on the test risk of
overparameterized regression involves a noise-propagation operator
$\bM = \bSigma^{-1}\bGamma\bSigma^{-1}$ built from the SVD of the
training Jacobian $\bJ$ and the test--train cross-correlation
$\bGamma$ of gradient features. The trichotomy
$\tr(\bM) < \infty \iff \beta > \alpha + 1/2$, under power-law
decay $\sigma_{j} \asymp j^{-\alpha}$ and
$\varrho_{j} := \sqrt{\bGamma_{jj}} \asymp j^{-\beta}$, characterizes
benign, tempered, and catastrophic regimes for squared loss; the
power-law model is standard in the random-feature analysis
of~\cite{mei2022}, with bias--variance refinements
in the companion~\cite{paperB}.

The same generalization story should apply to classification, but
the squared-loss derivation does not directly transfer: cross-entropy
introduces a nonlinear softmax link, and its parameter Hessian is
$\bJ^{\top}\bH\bJ$ rather than $\bJ^{\top}\bJ$. The natural question
is whether a Fisher-corrected operator preserves the trichotomy and
admits an exact effective-rank calculation at the Neural Collapse
(NC) fixed point of~\cite{papyan2020}.

This paper answers both questions affirmatively. We define
$\bM_{\mathrm{CE}}$ as the natural generalization of $\bM$ obtained
by replacing $\bJ^{\top}\bJ$ with the cross-entropy Hessian (see
Definition~\ref{def:M-CE} and the derivation that precedes it).
We then show:

\begin{enumerate}[label=(\roman*)]
\item The CE trichotomy
$\tr(\bM_{\mathrm{CE}}) < \infty \iff \beta > \alpha + 1/2$
holds whenever the predicted probabilities are bounded away from the
one-hot vertices uniformly along the relevant directions
(Theorem~\ref{thm:ce-trichotomy}).

\item Approaching the NC simplex equiangular tight frame fixed
point, each per-sample block of $\bM_{\mathrm{CE}}$ on the zero-sum
subspace has effective rank exactly $K-1$ (Theorem~\ref{thm:ce-rank}),
where $K$ is the number of classes; the same holds for the averaged
operator $\overline{\bM}_{\mathrm{CE}}$.
\end{enumerate}

The first result extends the trichotomy across loss functions; the
second states that, at the terminal phase of training, the
generalization-controlling rank is set by the number of classes
minus one, regardless of the input dimension. We work conditional
on the presence of NC and characterize what this implies for the
effective rank, rather than proving NC.

\paragraph{Relation to prior work.}
The Neural Collapse phenomenon was identified empirically
in~\cite{papyan2020} and characterized mathematically in
the unconstrained-features setting under cross-entropy
by~\cite{lu2022, mixon2022} and under squared loss / central path
by~\cite{han2022}. The class-imbalanced extension of unconstrained
features is~\cite{Fang2021}. The cross-entropy Fisher matrix
$\bH = \diag(\bp) - \bp\bp^{\top}$ and its degeneracy at one-hot
vertices is classical (see~\cite{Amari2016}); we use only the
finite-rank facts about $\bH$ and the simplex ETF geometry.

\paragraph{Organization.}
Section~\ref{sec:setting} fixes notation and recalls the
classification setting. Section~\ref{sec:operator} derives
$\bM_{\mathrm{CE}}$ from the CE Hessian. Section~\ref{sec:trichotomy}
proves the CE trichotomy. Section~\ref{sec:nc} derives the
effective-rank formula at the NC fixed point. Section~\ref{sec:disc}
discusses limitations.

%% ============================================================
\section{Setting and notation}\label{sec:setting}
%% ============================================================

Let $K \geq 2$ be the number of classes. The predictor has
$K$-dimensional output
$f(\bx;\theta) = (f_{1}(\bx;\theta),\ldots,f_{K}(\bx;\theta)) \in \R^{K}$,
softmax probabilities
$p_{k}(\bx;\theta) = e^{f_{k}}/\sum_{k'} e^{f_{k'}}$, and one-hot
labels $\by_{i} \in \{e_{1},\ldots,e_{K}\}$, where $e_{k}$ is the
$k$-th standard basis vector. Training data
$S = \{(\bx_{i},\by_{i})\}_{i=1}^{n}$ are i.i.d. The cross-entropy
loss is
\[
  \ell_{\mathrm{CE}}(f,\by) = -\sum_{k=1}^{K} y_{k}\log p_{k}(f).
\]
A direct computation gives the well-known output-space Hessian
\begin{equation}\label{eq:Hbox}
  \nabla_{f}^{2}\ell_{\mathrm{CE}}(f,\by) = \bH(\bx) := \diag(\bp(\bx)) - \bp(\bx)\bp(\bx)^{\top},
\end{equation}
which is symmetric positive semidefinite, has rank $K-1$ for any
$\bp$ in the open simplex (it has a single zero eigenvalue along
$\mathbf{1}$), and rank $0$ at every one-hot vertex
$\bp = e_{k}$ (since at $\bp = e_{k}$ we have
$\diag(e_{k}) = e_{k}e_{k}^{\top}$, hence
$\bH = \diag(e_{k}) - e_{k}e_{k}^{\top} = 0$).
The Fisher matrix $\bH(\bx)$ is the building block of the parameter
Hessian:
\begin{equation}\label{eq:CE-Hessian}
  \nabla_{\theta}^{2}\ell_{\mathrm{CE}}(\bx,\by;\theta)
  = \bJ(\bx)^{\top}\bH(\bx)\bJ(\bx) + \text{(label-dependent rank-1 term)},
\end{equation}
where $\bJ(\bx) \in \R^{K\times p}$ is the per-sample Jacobian of
$f(\bx;\cdot)$ with respect to $\theta$. The label-dependent term
vanishes in expectation over labels at any predictor that is exact
on the population (or at any limit point with vanishing gradient),
and we drop it from this point onward.

We adopt the notation of~\cite{paperA}. Let
$\bJ_{\mathrm{CE}} \in \R^{nK \times p}$ denote the stacked
training Jacobian with rows
$\nabla_{\theta} f_{k}(\bx_{i};\theta^{*})^{\top}$ ranging over
$(i,k) \in [n]\times[K]$, and let
$\bH = \bigoplus_{i=1}^{n}\bH(\bx_{i})$ denote the block-diagonal
sample Fisher. Denote the cross-class test--train cross-correlation
\begin{equation}\label{eq:GammaCE}
  (\bGamma_{\mathrm{CE}})_{(j,k),(j',k')} = \E_{\bx}[\psi_{j,k}(\bx)\psi_{j',k'}(\bx)],
  \quad \psi_{j,k}(\bx) = \nabla_{\theta} f_{k}(\bx;\theta^{*})^{\top}\bv_{j},
\end{equation}
where $\{\bv_{j}\}$ is the right SVD basis of an appropriate factor
of $\bJ_{\mathrm{CE}}$ (see Section~\ref{sec:operator}).

%% ============================================================
\section{Deriving $\bM_{\mathrm{CE}}$ from the cross-entropy
  Hessian}\label{sec:operator}
%% ============================================================

In MSE, the operator $\bM = \bSigma^{-1}\bGamma\bSigma^{-1}$ arises
because the test risk decomposes as a quadratic form in the
predicted residuals, weighted by the inverse of $\bJ^{\top}\bJ$ at
the interpolant. Concretely, $\bSigma$ comes from the SVD of the
training Jacobian, and the noise-leakage variance equals
$\sigma^{2}\tr(\bSigma^{-1}\bGamma\bSigma^{-1})$.

For cross-entropy at a (near-)interpolating predictor, the
linearized test risk is governed by the inverse of the parameter
Hessian~\eqref{eq:CE-Hessian}, that is
$(\bJ_{\mathrm{CE}}^{\top}\bH\bJ_{\mathrm{CE}})^{-1}$. Two natural
factorizations arise:

\begin{enumerate}
\item Define the \emph{Fisher-weighted Jacobian}
\[
  \widetilde\bJ := \bH^{1/2}\bJ_{\mathrm{CE}} \in \R^{nK\times p},
\]
with SVD $\widetilde\bJ = \widetilde\bU\,\widetilde\bSigma\,\widetilde\bV^{\top}$.
Then
$\bJ_{\mathrm{CE}}^{\top}\bH\bJ_{\mathrm{CE}} = \widetilde\bV\,\widetilde\bSigma^{2}\,\widetilde\bV^{\top}$,
so the parameter Hessian's eigenstructure coincides with that of
$\widetilde\bJ^{\top}\widetilde\bJ$.

\item The right singular vectors $\widetilde\bV$ provide the natural
basis for gradient features
$\widetilde\psi_{j,k}(\bx) := \nabla_{\theta}f_{k}(\bx;\theta^{*})^{\top}\widetilde\bv_{j}$,
and the corresponding cross-correlation matrix is
$\widetilde\bGamma_{(j,k),(j',k')} = \E_{\bx}[\widetilde\psi_{j,k}(\bx)\widetilde\psi_{j',k'}(\bx)]$.
The matrix $\widetilde\bSigma_{\mathrm{CE}}$ refers throughout to
the singular value matrix $\widetilde\bSigma$ above; these are the
square roots of the eigenvalues of the parameter Hessian.
\end{enumerate}

These two ingredients give the operator that propagates label noise
through the linearized CE risk:

\begin{definition}[Cross-entropy noise propagation operator]\label{def:M-CE}
The \emph{cross-entropy noise propagation operator} is
\begin{equation}\label{eq:M-CE}
  \bM_{\mathrm{CE}}
  := \widetilde\bSigma^{-1}\,\widetilde\bGamma\,\widetilde\bSigma^{-1},
  \qquad \widetilde\bJ = \bH^{1/2}\bJ_{\mathrm{CE}} = \widetilde\bU\widetilde\bSigma\widetilde\bV^{\top},
\end{equation}
where the inverse on the right is the Moore--Penrose pseudoinverse
restricted to the range of $\widetilde\bSigma$.
\end{definition}

\begin{definition}[Effective rank]\label{def:reff}
For any non-zero positive semidefinite operator $\bA$ on a finite-dimensional
inner-product space, the \emph{effective rank} is
\[
  r_{\mathrm{eff}}(\bA) \;:=\; \frac{(\tr\bA)^{2}}{\norm{\bA}_{F}^{2}}.
\]
This quantity satisfies $1 \leq r_{\mathrm{eff}}(\bA) \leq \rank(\bA)$,
with equality on the right iff all nonzero eigenvalues are equal,
and is invariant under positive scalar rescaling of $\bA$.
\end{definition}

\begin{remark}[Equivalence with the form
$\bF^{-1}\bGamma\bF^{-1}$]\label{rem:M-CE-eq}
Some prior writeups state $\bM_{\mathrm{CE}} = \bF^{-1}\bGamma_{\mathrm{CE}}\bF^{-1}$
with $\bF = \bH^{1/2}\bSigma_{\mathrm{CE}}$, where $\bSigma_{\mathrm{CE}}$
denotes the singular values of $\bJ_{\mathrm{CE}}$. This is the same
operator as~\eqref{eq:M-CE} written in the basis of $\bJ_{\mathrm{CE}}$
rather than $\widetilde\bJ$: $\widetilde\bSigma$ are precisely the
singular values of $\bH^{1/2}\bJ_{\mathrm{CE}}$, which equal the
square roots of the eigenvalues of
$\bJ_{\mathrm{CE}}^{\top}\bH\bJ_{\mathrm{CE}}$. We retain
Definition~\ref{def:M-CE} as the canonical form because the
factorization is unambiguous.
\end{remark}

The justification for Definition~\ref{def:M-CE} is the following
linearization lemma, which derives the noise-leakage variance from
a Taylor expansion of the empirical CE risk near $\theta^{*}$. The
argument is the cross-entropy analog of the MSE template
in~\cite[\S 2]{paperA}; the loss-specific ingredient is the
substitution of the CE Hessian $\bJ^{\top}\bH\bJ$ for the squared-loss
Hessian $\bJ^{\top}\bJ$, equivalently the replacement of $\bJ$ by
$\widetilde\bJ = \bH^{1/2}\bJ$.

\begin{lemma}[Linearized noise leakage under CE]\label{lem:noise-leakage}
Let $\theta^{*}$ be a stationary point of the empirical CE risk on
the training set $S$, with parameter Hessian
$\bF := \bJ_{\mathrm{CE}}^{\top}\bH\bJ_{\mathrm{CE}}$ invertible
on its range. Suppose label perturbations
$\bxi_{i} \in \R^{K}$ satisfy $\E[\bxi_{i}] = 0$ and
$\E[\bxi_{i}\bxi_{i}^{\top}] = \sigma^{2}\,\bH(\bx_{i})$ (the
Fisher-weighted noise model that matches the CE link). Let
$\theta^{*}_{\bxi}$ denote the perturbed stationary point. Then to
linear order in $\bxi$,
\begin{enumerate}[label=(\roman*)]
\item the parameter perturbation is
$\delta\theta = \bF^{\dagger}\bJ_{\mathrm{CE}}^{\top}\bH^{1/2}\bxi
              = \widetilde\bV\widetilde\bSigma^{-1}\widetilde\bU^{\top}\bxi$;
\item along any gradient feature direction $\widetilde\bv_{j}$
the variance contribution to the test residual at $\bx_{*}$ equals
$\sigma^{2}\widetilde\bGamma_{jj}/\widetilde\sigma_{j}^{2}$;
\item summing,
$\Var(\text{linearized test residual}) = \sigma^{2}\tr(\bM_{\mathrm{CE}})$.
\end{enumerate}
\end{lemma}

\begin{proof}
Linearize the empirical-risk gradient near $\theta^{*}$:
\[
  \nabla_{\theta}\widehat{\mathcal{R}}_{\mathrm{CE}}(\theta^{*} + \delta\theta;\, \by + \bxi)
   = \bF\,\delta\theta - \bJ_{\mathrm{CE}}^{\top}\bH^{1/2}\bxi
   + O(\norm{\delta\theta}^{2} + \norm{\bxi}^{2}),
\]
using $\nabla_{\theta}\ell_{\mathrm{CE}}(\bx,\by;\theta) = \bJ(\bx)^{\top}(\bp(\bx;\theta) - \by)$
and the noise model. Setting the gradient to zero gives
$\delta\theta = \bF^{\dagger}\bJ_{\mathrm{CE}}^{\top}\bH^{1/2}\bxi$,
which is~(i); substituting the SVD
$\widetilde\bJ = \widetilde\bU\widetilde\bSigma\widetilde\bV^{\top}$
into $\bF = \widetilde\bV\widetilde\bSigma^{2}\widetilde\bV^{\top}$
yields the explicit form. For (ii), the test residual contribution
along $\widetilde\bv_{j}$ is
$\widetilde\psi_{j,\cdot}(\bx_{*})^{\top}\delta\theta$, whose
variance is
$\sigma^{2}\widetilde\sigma_{j}^{-2}\widetilde\bGamma_{jj}$ by
orthonormality of $\widetilde\bU$. For (iii), sum to recognize
$\sum_{j}\widetilde\bGamma_{jj}/\widetilde\sigma_{j}^{2}
   = \tr(\widetilde\bSigma^{-1}\widetilde\bGamma\widetilde\bSigma^{-1})
   = \tr(\bM_{\mathrm{CE}})$.
\end{proof}

This lemma plays the same role as the squared-loss derivation
in~\cite{paperA}: $\bM_{\mathrm{CE}}$ is exactly the quadratic form
that controls noise-leakage variance under CE, with $\bH$ replacing
the identity in the output-space inner product.

%% ============================================================
\section{Cross-entropy trichotomy}\label{sec:trichotomy}
%% ============================================================

\begin{theorem}[Cross-entropy trichotomy]\label{thm:ce-trichotomy}
Suppose the singular values $\widetilde\sigma_{j}$ of $\widetilde\bJ
= \bH^{1/2}\bJ_{\mathrm{CE}}$ and the diagonal entries
$\widetilde\bGamma_{jj}$ (in the right singular basis of
$\widetilde\bJ$) exhibit power-law decay
$\widetilde\sigma_{j} \asymp j^{-\alpha}$,
$\widetilde\varrho_{j} := \sqrt{\widetilde\bGamma_{jj}} \asymp j^{-\beta}$.
Suppose furthermore that the Fisher block is uniformly non-degenerate
in the sense of the operator inequality
\begin{equation}\label{eq:hmin-op}
  \bH(\bx_{i}) \;\succeq\; h_{\min}\,\bP_{\mathcal{Z}_{K}}
  \qquad \text{for all } i = 1,\ldots,n,
\end{equation}
for some $h_{\min} > 0$, where $\bP_{\mathcal{Z}_{K}}$ is the
orthogonal projection onto the zero-sum subspace
$\mathcal{Z}_{K} = \{u \in \R^{K} : \mathbf{1}^{\top}u = 0\}$.
Then
\[
  \tr(\bM_{\mathrm{CE}}) < \infty \iff \beta > \alpha + \tfrac{1}{2}.
\]
\end{theorem}

\begin{proof}
By Definition~\ref{def:M-CE},
\[
  \tr(\bM_{\mathrm{CE}})
  = \sum_{j}\widetilde\bGamma_{jj}/\widetilde\sigma_{j}^{2}
  \asymp \sum_{j} j^{-2\beta}/j^{-2\alpha}
  = \sum_{j} j^{2(\alpha-\beta)},
\]
which converges iff $2(\alpha-\beta) < -1$, i.e.\ $\beta > \alpha + 1/2$.
This is a statement about the singular values of
$\widetilde\bJ$ and the diagonal of $\widetilde\bGamma$ in
$\widetilde\bJ$'s right singular basis, so no basis conversion is
needed. Hypothesis~\eqref{eq:hmin-op} ensures that
$\widetilde\bJ = \bH^{1/2}\bJ_{\mathrm{CE}}$ has full rank on the
zero-sum subspace whenever $\bJ_{\mathrm{CE}}$ does, with smallest
nonzero singular value at least $h_{\min}^{1/2}$ times that of
$\bJ_{\mathrm{CE}}$ on $\mathcal{Z}_{K}$, so the power-law
hypotheses are equivalent to power laws on $\bJ_{\mathrm{CE}}$ up to
a multiplicative constant.
\end{proof}

\begin{remark}[On the uniform Fisher lower bound]\label{rem:hmin}
Hypothesis~\eqref{eq:hmin-op} is the cross-entropy analog of the
``well-conditioned predictor'' assumption: at any predictor whose
softmax outputs $\bp(\bx)$ avoid the boundary of the simplex, the
nonzero eigenvalues of $\bH(\bx)$ are bounded below by a positive
constant on $\mathcal{Z}_{K}$. For a stable training trajectory
away from the zero-margin classifier, this is automatic with
$h_{\min}$ controlled by the entropy of the outputs. The hypothesis
fails in the strict NC limit; we treat that limit separately in
Section~\ref{sec:nc} via a different argument.
\end{remark}

\begin{remark}\label{rem:Sigma-CE}
The matrix $\bSigma_{\mathrm{CE}}$ in the form
$\bM_{\mathrm{CE}} = \bF^{-1}\bGamma_{\mathrm{CE}}\bF^{-1}$ of
Remark~\ref{rem:M-CE-eq} denotes the diagonal matrix of singular
values of the (unweighted) Jacobian $\bJ_{\mathrm{CE}}$; under
hypothesis~\eqref{eq:hmin-op}, the singular values of
$\widetilde\bJ$ relate to those of $\bJ_{\mathrm{CE}}$ via
$\widetilde\sigma_{j} \in [h_{\min}^{1/2}\sigma_{j}^{\mathrm{CE}},\, h_{\max}^{1/2}\sigma_{j}^{\mathrm{CE}}]$
where $h_{\max} = \norm{\bH}_{\mathrm{op}}$.
\end{remark}

%% ============================================================
\section{Effective rank at the Neural Collapse fixed
  point}\label{sec:nc}
%% ============================================================

We now compute the effective rank
$r_{\mathrm{eff}}(\bM_{\mathrm{CE}})$ in the limit where the
trained predictor approaches the Neural Collapse fixed point
of~\cite{papyan2020}. A subtlety: in the strict infinite-margin
limit (logits diverging in the simplex ETF directions), the softmax
output for each sample becomes one-hot, $\bH(\bx) \to 0$, and
$\bM_{\mathrm{CE}}$ degenerates. We address this by considering a
sequence of predictors whose last-layer features approach the NC
fixed point with \emph{uniformly bounded logit norms}; in this
regime the limiting $\bp(\bx)$ is a probability vector strictly in
the open simplex, peaked on the correct class but with positive
mass on every coordinate. This is the regime relevant to NC
phenomenology at any finite training time, and the one in which the
unconstrained-features analyses of~\cite{lu2022,mixon2022,Fang2021}
yield well-defined limits. The strict-margin limit is treated
separately by rescaling (Remark~\ref{rem:strict-NC}).

\subsection{Geometry of the simplex equiangular tight frame}

Recall the NC fixed point: the last-layer features satisfy
$h(\bx_{i}) = \mu_{c(\bx_{i})}$ where $\{\mu_{c}\}_{c=1}^{K}$
forms a simplex equiangular tight frame in $\R^{K-1}$ with
$\sum_{c}\mu_{c} = 0$ and
$\mu_{c}^{\top}\mu_{c'} = \tfrac{K\delta_{cc'} - 1}{K-1}\,\norm{\mu}^{2}$
for some common norm $\norm{\mu}$. The classifier weights
$\{\bw_{c}\}$ likewise satisfy
$\bw_{c} = (\norm{\bw}/\norm{\mu})\mu_{c}$ up to a scalar.

The image of the per-class logit difference map
$h \mapsto (\bw_{1}^{\top}h,\ldots,\bw_{K}^{\top}h)$ at the NC fixed
point therefore lies in the $(K-1)$-dimensional zero-sum subspace
\[
  \mathcal{Z}_{K} := \{u \in \R^{K} : \mathbf{1}^{\top}u = 0\}.
\]
Let $\bP_{\mathcal{Z}}$ be the orthogonal projection onto
$\mathcal{Z}_{K}$, that is
$\bP_{\mathcal{Z}} = \bI_{K} - \tfrac{1}{K}\mathbf{1}\mathbf{1}^{\top}$.

\subsection{Limiting Fisher and limiting cross-correlation}

\begin{lemma}[Fisher matrix under bounded-logit NC]\label{lem:H-NC}
Consider a sequence of predictors $\{\theta^{(t)}\}$ whose
last-layer features approach the NC fixed point and whose logit
vectors $f(\bx_{i};\theta^{(t)}) \in \R^{K}$ satisfy
$\norm{f(\bx_{i};\theta^{(t)})}_{\infty} \leq L$ for some finite
$L > 0$. Then for any sample $\bx_{i}$ and along any limit point of
$\{\theta^{(t)}\}$, every coordinate of the predicted probability
$\bp(\bx_{i};\theta^{(t)})$ is bounded below by
$p_{\min} := e^{-2L}/K > 0$, the matrix $\bH(\bx_{i})$ has rank
exactly $K-1$ with kernel $\spn(\mathbf{1})$, and the operator
inequality
\[
  \bH(\bx_{i}) \;\succeq\; h_{\min}\,\bP_{\mathcal{Z}_{K}},
  \qquad h_{\min} := p_{\min}(1 - 1/K),
\]
holds uniformly in $i$ and $t$.
\end{lemma}

\begin{proof}
For any logit vector $f$ with $\norm{f}_{\infty} \leq L$, the
softmax probabilities satisfy $p_{k} = e^{f_{k}}/\sum_{k'}e^{f_{k'}} \geq e^{-L}/(Ke^{L}) = e^{-2L}/K$,
giving $p_{\min} > 0$. The Fisher matrix
$\bH = \diag(\bp) - \bp\bp^{\top}$ satisfies
$\bH\mathbf{1} = \bp - \bp(\mathbf{1}^{\top}\bp) = 0$, so
$\spn(\mathbf{1}) \subseteq \ker\bH$, with equality on the open
simplex (a standard fact about the multinomial Fisher matrix; see
e.g.~\cite[\S 3]{Amari2016}). The bound on the smallest nonzero
eigenvalue of $\bH$ on $\mathcal{Z}_{K}$ follows from the standard
estimate $\lambda_{\min}(\bH|_{\mathcal{Z}_{K}}) \geq p_{\min}(1 - p_{\max})
\geq p_{\min}(1 - 1/K)$ for $\bp$ supported on the open simplex
with $p_{\max} \leq 1 - (K-1)p_{\min} \leq 1 - 1/K$ (achieved at
the corner of the bounded-logit constraint set).
\end{proof}

\begin{remark}[Strict-one-hot limit and rescaling]\label{rem:strict-NC}
In the strict infinite-margin limit ($L \to \infty$), the
bounded-logit condition fails, $\bp \to e_{c}$, and
$\bH \to 0$, so $\widetilde\bJ \to 0$ and Definition~\ref{def:M-CE}
degenerates. The natural rescaled object is
\[
  \widetilde\bJ_{\mathrm{rsc}} := h_{\min}^{-1/2}\widetilde\bJ
   \;=\; (h_{\min}^{-1/2}\bH^{1/2})\,\bJ_{\mathrm{CE}},
\]
which converges to the projector $\bP_{\mathcal{Z}_{K}}\bJ_{\mathrm{CE}}$
in the limit (since $h_{\min}^{-1/2}\bH^{1/2}$ has unit spectrum on
the support of $\bH$ and zero kernel along $\mathbf{1}$). The
operator $\bM_{\mathrm{CE}}$ rebuilt from
$\widetilde\bJ_{\mathrm{rsc}}$ is well-defined on the zero-sum
subspace and satisfies
$r_{\mathrm{eff}}(\bM_{\mathrm{CE}}^{\mathrm{rsc}}) = r_{\mathrm{eff}}(\bM_{\mathrm{CE}})$
since effective rank is invariant under positive scalar rescaling
of $\widetilde\bSigma$. In what follows we work with finite logit
scale where this rescaling is unnecessary.
\end{remark}

\subsection{Effective rank formula}

Before stating the theorem, we make the structural NC hypotheses
explicit. The simplex-ETF geometry of~\cite{papyan2020} fixes the
last-layer features and classifier weights, but the per-sample
Jacobian also depends on the backbone (everything below the last
layer). The following assumption, standard in the
unconstrained-features literature
\cite{lu2022,mixon2022,Fang2021}, isolates what is needed.

\begin{assumption}[NC backbone non-degeneracy]\label{ass:nc-backbone}
At $\theta^{*}$, with last-layer feature map
$h(\bx;\theta^{*}) \in \R^{K-1}$ and last-layer classifier
$W \in \R^{K\times(K-1)}$ in simplex-ETF configuration:
\begin{enumerate}[label=(\alph*)]
\item (Class collapse / NC1) For every sample $\bx_{i}$ in class
$c_{i} \in \{1,\ldots,K\}$, $h(\bx_{i};\theta^{*}) = \mu_{c_{i}}$
where $\{\mu_{c}\}$ is the simplex ETF
($\sum_{c}\mu_{c} = 0$,
$\mu_{c}^{\top}\mu_{c'} = (K\delta_{cc'} - 1)/(K-1) \cdot \norm{\mu}^{2}$).
\item (Self-duality / NC2) $W = (\norm{W}/\norm{\mu})\,M$ where
$M^{\top} = [\mu_{1}\,\cdots\,\mu_{K}]$.
\item (Backbone full-rank in ETF directions) The backbone Jacobian
$\partial h(\bx_{i};\theta^{*})/\partial\theta \in \R^{(K-1) \times p}$
has full row rank $K-1$ for every $i$.
\end{enumerate}
\end{assumption}

Hypothesis (c) is the chain-rule input to the per-sample Jacobian
$\bJ(\bx_{i}) \in \R^{K\times p}$; it follows in the
unconstrained-features regime where $h$ is a free variable
\cite{mixon2022}, and is the operative assumption in feature-learning
analyses of NC. We do not derive it from NC1--NC4; we assume it.

We measure the effective rank on the per-sample block (equivalently,
on the output-averaged operator), which is the geometric object
fixed by simplex-ETF symmetry and is independent of $n$.

\begin{definition}[Per-sample / averaged $\bM_{\mathrm{CE}}$ block]\label{def:M-bar}
Let $\bM_{\mathrm{CE}}^{(i)} \in \R^{K\times K}$ denote the
per-sample restriction of $\bM_{\mathrm{CE}}$ to the output block
indexed by sample $i$, viewed as an operator on $\R^{K}$. Define
the \emph{averaged operator}
\[
  \overline{\bM}_{\mathrm{CE}}
  := \frac{1}{n}\sum_{i=1}^{n}\bM_{\mathrm{CE}}^{(i)}.
\]
\end{definition}

\begin{theorem}[Effective rank at NC]\label{thm:ce-rank}
Let $\theta^{*}$ be a limit point of the training trajectory of
Lemma~\ref{lem:H-NC} satisfying Assumption~\ref{ass:nc-backbone},
and assume the training set is class-balanced
($n/K$ samples per class). Then both per-sample blocks and the
averaged operator satisfy
\[
  r_{\mathrm{eff}}(\bM_{\mathrm{CE}}^{(i)}) \;=\; K - 1
  \qquad \text{and} \qquad
  r_{\mathrm{eff}}(\overline{\bM}_{\mathrm{CE}}) \;=\; K - 1.
\]
\end{theorem}

\begin{proof}
\emph{Step 1: Per-sample block has rank exactly $K-1$ on
$\mathcal{Z}_{K}$.}
By Lemma~\ref{lem:H-NC}, $\bH(\bx_{i})$ is positive definite on
$\mathcal{Z}_{K}$ and zero on $\spn(\mathbf{1})$. By
Assumption~\ref{ass:nc-backbone}(c) the backbone Jacobian has rank
$K-1$, and the chain rule gives
\[
  \bJ(\bx_{i}) \;=\; W \cdot \frac{\partial h(\bx_{i};\theta^{*})}{\partial\theta}
                \;+\; \frac{\partial f}{\partial W}\Big|_{(\bx_{i})} \cdot
                  \frac{\partial W}{\partial\theta},
\]
where $\partial f/\partial W$ at the NC weight has rank $K-1$ on
$\mathcal{Z}_{K}$ (the simplex-ETF rows of $W$ span $\mathcal{Z}_{K}$
exactly, by Assumption~\ref{ass:nc-backbone}(b)). Hence
$\bJ(\bx_{i})\bP_{\mathcal{Z}_{K}}$ has row rank exactly $K-1$, and
the same holds for $\widetilde\bJ(\bx_{i}) = \bH(\bx_{i})^{1/2}\bJ(\bx_{i})$.

\emph{Step 2: Per-sample block has all $K-1$ eigenvalues equal.}
By Lemma~\ref{lem:H-NC} and the simplex-ETF symmetry, the per-sample
Fisher restricted to $\mathcal{Z}_{K}$ is a positive scalar multiple
of the identity:
$\bH(\bx_{i})|_{\mathcal{Z}_{K}} = h_{i}\,\bI_{\mathcal{Z}_{K}}$ for
some $h_{i} > 0$. (Direct calculation: at NC, $\bp(\bx_{i})$ is
permutation-symmetric across the $K-1$ non-target classes by the
simplex-ETF symmetry, so $\bH(\bx_{i})|_{\mathcal{Z}_{K}}$ is a
scalar matrix.) Likewise, by simplex-ETF symmetry across the $K-1$
output directions in $\mathcal{Z}_{K}$, the cross-correlation
$\widetilde\bGamma$ restricted to a per-sample block on
$\mathcal{Z}_{K}$ is a scalar matrix
$\gamma_{i}\,\bI_{\mathcal{Z}_{K}}$ for some $\gamma_{i} \geq 0$:
the $K-1$ within-block directions are equivalent under the action
of the simplex-ETF symmetry group $S_{K-1}$ that fixes the class of
$\bx_{i}$.

\emph{Step 3: Computing $r_{\mathrm{eff}}$.}
By Steps 1--2, $\bM_{\mathrm{CE}}^{(i)}$ on $\mathcal{Z}_{K}$ is
$(\gamma_{i}/h_{i})\bI_{\mathcal{Z}_{K}}$ — a positive scalar
multiple of the identity on a $(K-1)$-dimensional space, with the
remaining direction $\spn(\mathbf{1})$ killed by $\bH$. All $K-1$
nonzero eigenvalues of $\bM_{\mathrm{CE}}^{(i)}$ are equal, hence by
Definition~\ref{def:reff},
\[
  r_{\mathrm{eff}}(\bM_{\mathrm{CE}}^{(i)})
  = \frac{((K-1)\,\gamma_{i}/h_{i})^{2}}{(K-1)\,(\gamma_{i}/h_{i})^{2}}
  = K - 1.
\]

\emph{Step 4: Averaged operator.}
By class-balance, the per-sample blocks group into $K$ equal-size
class subsets, each of which gives the same $\gamma/h$ ratio by
the simplex-ETF symmetry permuting classes. Hence
$\overline{\bM}_{\mathrm{CE}} = (\bar\gamma/\bar h)\bI_{\mathcal{Z}_{K}}$
on $\mathcal{Z}_{K}$ for the corresponding averages, with all $K-1$
nonzero eigenvalues equal. By the same calculation,
$r_{\mathrm{eff}}(\overline{\bM}_{\mathrm{CE}}) = K-1$.
\end{proof}

\begin{remark}[Dimension-free scaling for
classification]\label{rem:dim-free}
Theorem~\ref{thm:ce-rank} states that, at the NC fixed point, the
effective rank of $\bM_{\mathrm{CE}}$ on each per-sample block (and
of the averaged operator) depends only on $K$ and not on the input
dimension $d$ or the parameter count $p$. This is the classification
analog of the dimension-free regression rate (rank-one collapse
onto a single feature direction): the rank-$(K-1)$ collapse follows
from the geometric fact that the simplex ETF in $\R^{K-1}$ spans a
$(K-1)$-dimensional subspace, on which the Fisher matrix at NC is a
scalar.
\end{remark}

\begin{remark}[Full-operator effective rank]\label{rem:full-vs-per-sample}
The full operator $\bM_{\mathrm{CE}}$ on the joint sample-class
space has rank $n(K-1)$ on the zero-sum subspace by Step~1,
all eigenvalues equal under class-balance and simplex-ETF symmetry,
so $r_{\mathrm{eff}}(\bM_{\mathrm{CE}}) = n(K-1)$ at the limit
point. This is consistent: each of the $n$ training samples
contributes $K-1$ equal eigenvalues. The dimension-free scaling we
care about for generalization is the per-sample / averaged
quantity, which captures the noise leakage \emph{per training
example}.
\end{remark}

%% ============================================================
\section{Discussion}\label{sec:disc}
%% ============================================================

\paragraph{Status of the NC hypothesis.}
The conclusion $r_{\mathrm{eff}}(\bM_{\mathrm{CE}}^{(i)}) = K-1$
(per-sample) is conditional on the NC fixed point being attained or
approached at the end of training. NC is observed empirically
across diverse networks~\cite{papyan2020}, and is rigorously
established in the unconstrained-features model
of~\cite{lu2022,mixon2022}, which idealizes the last-layer features
as free variables. The result of this paper is therefore a
statement of the form ``NC $+$ backbone non-degeneracy
(Assumption~\ref{ass:nc-backbone}(c)) $\Rightarrow$ rank-$(K-1)$
effective rank''; we do not prove that NC holds for finite networks
trained end-to-end with cross-entropy. Limitations in the strict
one-hot limit are addressed by rescaling
(Remark~\ref{rem:strict-NC}); see also the discussion below.

\paragraph{Off-fixed-point behavior.}
A natural follow-up question is: how does
$r_{\mathrm{eff}}(\bM_{\mathrm{CE}})$ behave during the approach to
the NC fixed point? We expect a gap term whose decay is governed by
the convergence rate of NC, which would yield a rate of
$r_{\mathrm{eff}}(\bM_{\mathrm{CE}})(t) - (K-1)$ in terms of the
training time $t$. Quantifying this gap is open.

\paragraph{Class-imbalanced data.}
Theorem~\ref{thm:ce-rank} assumes class-balanced gradient features.
For imbalanced classes, the simplex ETF deforms to a
class-frequency-weighted variant (see e.g.~\cite{Fang2021} for the
unconstrained-features analysis), and the effective rank gains a
class-frequency-dependent correction. We do not pursue this here.

\paragraph{Connection to the regression result of~\cite{paperC}.}
The companion paper~\cite{paperC} (this paper's regression
counterpart) shows that, after feature-direction alignment, the
effective rank of the regression operator $\bM$ (the squared-loss
analog of $\bM_{\mathrm{CE}}$ in this paper) collapses to a single
dimension along the target direction. In classification, the
analogous collapse for $\bM_{\mathrm{CE}}$ is to $K-1$ dimensions
along the simplex faces: each class identifies its corresponding
simplex vertex, and the rank-one regression collapse for $\bM$
becomes a rank-$(K-1)$ classification collapse for
$\bM_{\mathrm{CE}}$. The two results share the same mathematical
mechanism: a geometric symmetry in the trained predictor's
last-layer geometry collapses an a priori $\Theta(d)$ rank to a
small intrinsic dimension. We retain distinct notation
$\bM$ vs.\ $\bM_{\mathrm{CE}}$ throughout to disambiguate the two
operators.

\subsection*{Acknowledgements}

The exploratory analysis underlying Theorems~\ref{thm:ce-trichotomy}
and~\ref{thm:ce-rank}, including the derivation of the
Fisher-corrected operator from the cross-entropy Hessian, was
conducted with the assistance of Claude (Anthropic). All
mathematical statements and proofs have been verified by the
author.

%% ============================================================
%% References
%% ============================================================

\begin{thebibliography}{20}

\bibitem{paperA}
L.~Chang.
\newblock Noise propagation operators and a two-parameter characterization
  of benign overfitting.
\newblock \emph{Companion paper}, 2026.

\bibitem{paperB}
L.~Chang.
\newblock A bias--variance decomposition for the noise propagation operator
  and the role of feature learning.
\newblock \emph{Companion paper}, 2026.

\bibitem{paperC}
L.~Chang.
\newblock Feature learning, effective dimension, and architecture-specific
  critical depth.
\newblock \emph{Companion paper}, 2026.

\bibitem{Amari2016}
S.-i.~Amari.
\newblock \emph{Information Geometry and Its Applications}.
\newblock Springer, 2016.

\bibitem{Fang2021}
C.~Fang, H.~He, Q.~Long, and W.~J.~Su.
\newblock Exploring deep neural networks via layer-peeled model:
  Minority collapse in imbalanced training.
\newblock \emph{Proceedings of the National Academy of Sciences},
  118(43):e2103091118, 2021.

\bibitem{han2022}
X.~Y.~Han, V.~Papyan, and D.~L.~Donoho.
\newblock Neural collapse under MSE loss: proximity to and dynamics on
  the central path.
\newblock In \emph{International Conference on Learning Representations
  (ICLR)}, 2022.

\bibitem{lu2022}
J.~Lu and S.~Steinerberger.
\newblock Neural collapse under cross-entropy loss.
\newblock \emph{Applied and Computational Harmonic Analysis},
  59:224--241, 2022.

\bibitem{mei2022}
S.~Mei, T.~Misiakiewicz, and A.~Montanari.
\newblock Generalization error of random feature and kernel methods:
  Hypercontractivity and kernel matrix concentration.
\newblock \emph{Applied and Computational Harmonic Analysis},
  59:3--84, 2022.

\bibitem{mixon2022}
D.~G.~Mixon, H.~Parshall, and J.~Pi.
\newblock Neural collapse with unconstrained features.
\newblock \emph{Sampling Theory, Signal Processing, and Data Analysis},
  20(2):11, 2022.

\bibitem{papyan2020}
V.~Papyan, X.~Y.~Han, and D.~L.~Donoho.
\newblock Prevalence of neural collapse during the terminal phase of deep
  learning training.
\newblock \emph{Proceedings of the National Academy of Sciences},
  117(40):24652--24663, 2020.

\end{thebibliography}

\end{document}
